\documentclass[12pt,a4paper]{article}
\usepackage[utf8]{inputenc}
\usepackage[portuguese]{babel}
\usepackage{amsmath}
\usepackage{amsfonts}
\usepackage{amssymb}
\usepackage{fancyhdr}
\usepackage{geometry}
\usepackage{graphicx}
\usepackage{tcolorbox}
\usepackage{parskip}
\usepackage{setspace}
\usepackage{tikz}
\usepackage{verbatim}
\usetikzlibrary{arrows.meta, calc}
\setstretch{1.2}
\tcbuselibrary{skins, breakable}
\geometry{top=2cm, bottom=2cm, left=2.5cm, right=2.5cm}
\usepackage[colorlinks=true, linkcolor=blue, urlcolor=blue, citecolor=blue]{hyperref}

% Estilo de caixa para destaque conceitual
\newtcolorbox{destaque}[1][]{
	enhanced,
	colback=gray!5,
	colframe=black!70,
	boxrule=1pt,
	arc=4pt,
	left=10pt,
	right=10pt,
	top=8pt,
	bottom=8pt,
	fonttitle=\bfseries,
	title=#1
}

\begin{document}

	
% --- CABEÇALHO COM CAIXA ---
\begin{tcolorbox}[
	enhanced, 
	colback=white, 
	colframe=black, 
	boxrule=2pt,
	left=0pt, 
	right=0pt,
	top=5pt,
	bottom=5pt
	]
	\noindent
	\begin{minipage}[c]{0.25\textwidth}
		\centering
		\includegraphics[width=4.5cm]{ifes.jpg} 
	\end{minipage}
	\begin{minipage}[c]{0.70\textwidth}
		\centering
		\textbf{Introdução ao Estudo das Funções -- 2026/2}\\
		\textbf{Licenciatura em Matemática}\\
		\textbf{Resumo de Função Quadrática}\\
		Professor: \textbf{Alex Jordane}
	\end{minipage}
\end{tcolorbox}

	\begin{destaque}[Forma Polinomial ou Geral]
		\label{dest:formageral}
		$$f(x) = ax^2 + bx + c, \quad \text{com } a, b, c \in \mathbb{R} \text{ e } a \neq 0$$
	\end{destaque}
	
	\begin{destaque}[Forma Canônica]
		\label{dest:formacanonica}
		$$f(x) = a(x - x_v)^2 + y_v$$
		onde $V(x_v, y_v)$ é o vértice da parábola, dado por:
		$$x_v = -\frac{b}{2a} \quad \text{e} \quad y_v = -\frac{\Delta}{4a}, \quad \text{com } \Delta = b^2 - 4ac$$
	\end{destaque}
	
	\begin{destaque}[Raízes da função]
		\label{dest:raizes}
		$$x_1  = \frac{-b - \sqrt{b^2 - 4ac}}{2a} \text{ e } x_2  = \frac{-b + \sqrt{b^2 - 4ac}}{2a}$$
	\end{destaque}		
	
	\begin{destaque}[Forma Fatorada]
		\label{dest:formafatorada}
		$$f(x) = a(x - x_1)(x - x_2), \quad \text{com } \Delta \geqslant 0$$
		onde $x_1$ e $x_2$ são as raízes reais da função ($f(x_1) = f(x_2) = 0$).
	\end{destaque}
	
	\begin{destaque}[Forma Soma-Produto]
		\label{dest:soma-produto}
			$$f(x) = a(x^2 - S \cdot x + P)$$
		onde $S = x_1 + x_2$ e $P=x_1 \cdot x_2$.
	\end{destaque}


\end{document}